\section{Техническое задание}
\subsection{Основание для разработки}

Основанием для разработки является задание на выпускную квалификационную работу бакалавра "<Разработка программной системы компилятора функционального языка программирования Common Lisp с оптимизациями">.

\subsection{Цель и назначение разработки}

Основной задачей выпускной квалификационной работы является разработка компилятора Common Lisp с оптимизациями и его тестирование.

Целью разработки является создание оптимизирующего компилятора Common Lisp.

Задачами данной разработки являются:
\begin{itemize}
\item изучение процесса компиляции языков программирования и функциональных языков в частности;
\item изучение методов оптимизации компиляторов;
\item проектирование компилятора с оптимизациями;
\item реализация компиляции S-выражения в промежуточную форму;
\item реализация генерации кода ассемблера;
\item реализация генерации байт-кода;
\item реализация оптимизации компиляции и генерации.
\end{itemize}

\subsection{Реализуемое подмножество Common Lisp}

\subsubsection{S-выражения}

Компилируемое выражение представляет собой S-выражение -- атом (символ, число, строка и др.) или список из нуля или более S-выражений.

\subsubsection{Типы объектов}

К объектам относятся числа (5, 0xf, 3.14), символы (AB), списки ( (1 2 3) ), пары ( (1 . 2) ), строки ("123"), массивы ( \#(1 2 3) ), одиночные символы (\#A), функции ( \#'(lambda (x) x) ). При переводе в логический тип NIL считается как ложь, всё остальное -- истина (при выводе истины используется символ T).

\subsubsection{Специальные формы}

\begin{enumerate}
\item (quote x) -- особый оператор, который возвращает свой аргумент, не вычисляя его. Сокращенная запись -- 'x (кавычка x).
\item (if <условие> <ветка true> <ветка false>). Условный оператор, считающий в своём условии истиной всё, кроме nil.
\item (setq <переменная$_1$> <выражение$_1$> ... <переменная$_n$> <выражение$_n$>). Создает глобальные переменные и присваивает им значения выражений. Возвращает значение последнего выражения.
\item (progn $e_1$ $e_2$ ... $e_n$). Последовательно вычисляет выражения $e_1$ ... $e_n$. Возвращает значение последнего выражения $e_n$.
\item (tagbody <список символов или выражений>). Последовательно вычисляет вложенные выражения. Символы выступают в роли меток, на которые можно перейти с помощью оператора go. Возвращает nil.
\item (catch <имя> <список выражений>). Блок с возможностью нелокального выхода. Сначала вычисляется имя блока, затем последовательно вычисляются выражения. Если встречается (throw <имя блока> <результат вычисления блока>), то вычисляется имя блока и происходит выход из ближайшего блока с этим именем. Если выход не происходит, то возвращается результат последнего выражения. Все имена вычисляются динамически.
\item (labels (<список функций>) <список выражений>). Объявляет локальные функции (в том числе и рекурсивные), после чего последовательно вычисляет выражения.
\end{enumerate}

\subsubsection{Примитивы}

\begin{itemize}

\paragraph{Списки}

\item (car x) -- возвращает первый элемент списка (левый объект пары);
\item (cdr x) -- возвращает все кроме первого элемента списка (правый объект пары);
\item (cons x y) -- создает точечную пару с объектами x и y;
\item (rplaca x y) -- заменяет CAR списка x на значение y;
\item (rplacd x y) -- заменяет CDR списка x на значение y;
\item (list $e_1$ $e_2$ ... $e_n$) -- создаёт список со значениями выражений $e_1$ $e_2$ ... $e_n$.

\paragraph{Арифметические операции}

\item (+ x y) -- сложение x и y;
\item (- x y) -- вычитание x и y;
\item (* x y) -- умножение x и y;
\item (/ x y) -- деление x и y;
\item (sin num) - синус вещественного числа num;
\item (cos num) - косинус вещественного числа num.

\paragraph{Предикаты}

\item (atom x) -- возвращает T (истина), если аргумент является атомом, иначе возвращает NIL, эквивалентный пустому списку (), что означает ложь;
\item (eq x y) -- возвращает T, если x и y -- один и тот же объект, иначе возвращает NIL;
\item (and <список условий>) -- логическая операция И;
\item (or <список условий>) -- логическая операция ИЛИ;
\item (equal x y) -- сравнение x и y;
\item (> x y) -- x больше чем y;
\item (< x y) -- x меньше чем y;
\item (\% x y) -- остаток деления x и y;
\item (pairp x) -- предикат, является ли x парой;
\item (symbolp x) -- предикат, является ли x символом;
\item (integerp x) -- предикат, является ли x числом;
\item (stringp x) -- предикат, является ли x строкой;
\item (arrayp x) -- предикат, является ли x массивом;
\item (floatp x) -- предикат, является ли x числом с плавающей запятой.

\paragraph{Побитовые операции}

\item (\& x y) -- побитовое И от x и y;
\item (bitor x y) -- побитовое ИЛИ от x и y;
\item (<{<} x y) -- сдвиг влево x на y;
\item (>{>} x y) -- сдвиг вправо x на y.

\paragraph{Строки и символы}

\item (intern str) -- создаёт символ на основе строки str;
\item (symbol-name sym) -- создаёт строку на основе символа sym;
\item (string-size str) -- возвращает длину строки str;
\item (concat $s_1$ ... $s_n$) -- объединяет строки в одну общую строку (возвращает пустую строку при отсутствии аргументов);
\item (inttostr x) -- перевод целочисленного числа x в строку;
\item (code-char x) -- создает символ по коду x;
\item (char-code c) -- возвращает код символа c;
\item (make-string len c) -- создаёт строку длиной len, заполненную символом c;
\item (char str i) -- получает символ строки str по индексу i;
\item (sets str i c) -- заменяет символ в строке str с индексом i символом c;
\item (subseq str a b) -- получение подстроки из строки str (от индекса a до b).

\paragraph{Массивы}

\item (make-array x) -- создание пустого массива длины x;
\item (array-size x) -- возвращает размер массива x;
\item (aref x i) -- чтение элемента массива x по индексу i;
\item (seta x i y) -- присвоение значения y элементу i массива x.

\paragraph{Печать объектов}

\item (print x) -- печать объекта x;
\item (putchar c) -- печать символа c.

\paragraph{Другие примитивы}

\item (funcall <функция> <аргументы>) -- вызывает функцию с аргументами;
\item (apply <функция> <список аргументов>) -- вызывает функцию с аргументами, переданными через список;
\item (eval <выражение>) -- вычисляет значение выражения как ещё одно выражение;
\item (error msg) -- останавливает вычисление, выводит сообщение об ошибке и возвращается в REPL цикл.

\end{itemize}

\subsubsection{Лямбда выражения}

Лямбда выражение -- это анонимная функция \#'(lambda ($p_1$ ... $p_n$) $e_1$ ... $e_n$), где $p_1$ ... $p_n$ -- это параметры функции, а $e_1$ ... $e_n$ -- выражения тела функции.

Вызов функции -- это следующее выражение:

\begin{figure}[H]
\begin{lstlisting}[language=Lisp]
(funcall #'(lambda (p1 ... pn) e1 ... en) a1 ... an)
\end{lstlisting}
\end{figure}

Сначала вычисляются все аргументы a1 ... an. Затем каждому параметру p1 ... pn ставится в соответствие вычисленное значение аргументов a1 ... an. После этого последовательно вычисляются выражения e1 ... en, содержащие параметры, вместо которых будут подставлены их значения.

\subsubsection{Определение функций}

Новую функцию можно создать с помощью оператора defun:

\begin{figure}[H]
\begin{lstlisting}[language=Lisp]
(defun null (x)
  (eq x ()))
\end{lstlisting}
\end{figure}

\subsubsection{Глобальные переменные}

Глобальные переменные существуют все время работы. Они создаются с помощью функции (defvar имя\_переменной [значение]). Значение может быть выражением:

\begin{figure}[H]
\begin{lstlisting}[language=Lisp]
(defvar a 10) -> A
A -> 10
\end{lstlisting}
\end{figure}

При отсутствии значения в переменную записывается значение NIL.

Установить значение переменной можно с помощью функции setq (особая форма). Если такой переменной не было то она создается:

\begin{figure}[H]
\begin{lstlisting}[language=Lisp]
(setq x 1) -> 1
x -> 1
\end{lstlisting}
\end{figure}

Можно одной функцией установить значения нескольких переменных:

\begin{figure}[H]
\begin{lstlisting}[language=Lisp]
  (setq a 1 b 2 c 3)
\end{lstlisting}
\end{figure}

Если переменная локальная (параметр функции), то setq ее модифицирует:

\begin{figure}[H]
\begin{lstlisting}[language=Lisp]
(defun test(x)
  (setq x 10)) ; модификация параметра
\end{lstlisting}
\end{figure}

Глобальные переменные являются динамическими, то есть их значение используется всегда последнее, которое было связано.

\subsubsection{Макросы}

Макрос задает шаблон для генерации выражения.

\begin{figure}[H]
\begin{lstlisting}[language=Lisp]
(defmacro test (var val)
  (list 'setq var val))
\end{lstlisting}
\end{figure}

При вызове макроса сначала происходит вычисление тела макроса (развертывание макроса):

\begin{figure}[H]
\begin{lstlisting}[language=Lisp]
(test abc 100) -> (setq abc 100)
\end{lstlisting}
\end{figure}

Затем получившееся выражение вычисляется:

\begin{figure}[H]
\begin{lstlisting}[language=Lisp]
(setq abc 100) ; abc = 100
\end{lstlisting}
\end{figure}

Обратная кавычка вычисляется как обычное цитирование, но также позволяет указывать, какие части цитирования должны быть вычислены. Эти части указываются с помощью запятой:

\begin{figure}[H]
\begin{lstlisting}[language=Lisp]
(setq a 10)
`(a b c ,a) -> (A B C 10)
\end{lstlisting}
\end{figure}

Запятая-at служит для того, чтобы подставить список (результат вычисления выражения внутри запятой-at должен быть списком):

\begin{figure}[H]
\begin{lstlisting}[language=Lisp]
(defvar a '(1 2 3)) -> (1 2 3)
`(,a ,@a) -> ((1 2 3) 1 2 3)
\end{lstlisting}
\end{figure}

Можно посмотреть результат макроподстановки с помощью функции macroexpand:

\begin{figure}[H]
\begin{lstlisting}[language=Lisp]
(macroexpand '(if (= 1 1) 2 3)) -> (COND ((= 1 1) 2) (T 3))
\end{lstlisting}
\end{figure}

\subsubsection{Функции как объекты первого класса}

Функции можно передавать в качестве параметров других функций и возвращать функции. Типы функций:

\begin{itemize}
\item локальные;
\item глобальные;
\item встроенные;
\item лямбда.
\end{itemize}

Пространство имен функций не совпадает с пространством имен переменных. Функция передается в качестве параметра как лямбда-выражение или символ. Но при вызове функции из параметра мы должны указать, что имя символа (переменная) относится к пространству имен функций. Для вызова функций по параметру используется примитив funcall. Пример:

\begin{figure}[H]
\begin{lstlisting}[language=Lisp]
; вызов примитива
(funcall #'+ 1 2 3) -> 6

; вызов lambda
(funcall #'(lambda (x y) (+ x y)) 2 3) -> 5

; вызов пользовательской функции по символу
(defun test (x) (+ x 1))
(funcall #'test 4) -> 5
\end{lstlisting}
\end{figure}

Оператор \#' является сокращением специальной формы function, создающей объект-замыкание, в котором сохраняется текущее окружение. При применении функции-замыкания будет использоваться не текущее окружение, а окружение, в котором было создано замыкание. Свободные локальные переменные являются лексическими, то есть сохраняют свое значение, которое было на момент замыкания. Глобальные переменные ведут себя динамическим образом - берется их значение в момент вызова функции.

\subsubsection{Обработка ошибок}

Обработка исключения для выражения:

\begin{figure}[H]
\begin{lstlisting}[language=Lisp]
(handle <выражение>
  (error (<параметры исключения>)
    <тело обработчика любых исключений в том числе просто error>)
  (division-by-zero (c)
    <тело обработчика деления на ноль>)
  ...
  (<пользовательское исключение> (<параметры>)
    <обработчик пользовательского исключения>))
\end{lstlisting}
\end{figure}

Для error существует обработчик по умолчанию, который выводит сообщение об ошибке. Примеры ошибок: division-by-zero (деление на ноль), index-error (выход за границы массива). Исключение может быть внутреннее или пользовательское:

\begin{figure}[H]
\begin{lstlisting}[language=Lisp]
(handle (if nil (raise 'my-error "My error" 20 30))
  (my-error (params)
  (print params)))
\end{lstlisting}
\end{figure}

\subsection{Программный интерфейс оптимизирующего компилятора}

Программный интерфейс компилятора должен состоять из единственной функции -- (compile expr optimization-flags), которая принимает на вход исходное S-выражение и флаги оптимизации optimization-flags и преобразует его в эквивалентный байт-код для выполнения на виртуальной машине.

Флаги оптимизации compiler-optimizations должны быть в виде списка символов, каждый из которых включает ту или иную оптимизацию компиляции. Пустой список выключает все оптимизации, специальный символ all -- включает все доступные оптимизации.

Результатом функции является список, первый элемент которого -- массив инструкций для целевой виртуальной машины, второй -- список констант, используемых в программе, и третий -- размер памяти для глобальных переменных. В случае ошибки компиляции возвращается строка с описанием ошибки.

\subsubsection{Флаги оптимизации компилятора}

Функция compile должна принимать следующие флаги оптимизации:

\begin{enumerate}
\item accumulator-loading -- удаление подряд идущих обращений к константам или переменным, которые не повлияют на результат вычисления выражения. Пример:
\begin{figure}[H]
\begin{lstlisting}[language=Lisp]
(let ((a 5)) 1 a 2) -> (let ((a 5)) 2)
\end{lstlisting}
\end{figure}

\item trivial-condition -- оптимизация тривиальных условий. Пример:
\begin{figure}[H]
\begin{lstlisting}[language=Lisp]
(if t 1 2) -> 1
\end{lstlisting}
\end{figure}

\item simplify-arithmetic -- оптимизация арифметических выражений (замена примитивов с произвольным числом аргументов на примитивы с двумя аргументами). Пример:
\begin{figure}[H]
\begin{lstlisting}[language=Lisp]
(+ 1 2 3 4) -> (+ (+ (+ 1 2) 3) 4)
\end{lstlisting}
\end{figure}

\end{enumerate}

%\begin{enumerate}
%\item constant-folding -- оптимизация избыточных вычислений с константами. Пример:
%\begin{figure}[H]
%\begin{lstlisting}[language=Lisp]
%(+ 1 2 3) -> 6
%\end{lstlisting}
%\end{figure}
%
%\item dead-code-elimination -- удаление выражений, которые не влияют на результат выполнения программы. Пример:
%\begin{figure}[H]
%\begin{lstlisting}[language=Lisp]
%(progn 1 2 3) -> (progn 3)
%\end{lstlisting}
%\end{figure}
%
%\item unreachable-code-elimination -- удаление недостижимого кода. Пример:
%\begin{figure}[H]
%\begin{lstlisting}[language=Lisp]
%(if t 1 2) -> 1
%\end{lstlisting}
%\end{figure}
%
%\item inlining -- встраивание функций. Пример:
%\begin{figure}[H]
%\begin{lstlisting}[language=Lisp]
%(progn
%  (defun f () 5)
%  (+ (f) 5))
%->
%(progn
%  (+ 5 5))
%\end{lstlisting}
%\end{figure}
%
%\item tail-call -- оптимизация хвостовой рекурсии. Пример:
%\begin{figure}[H]
%\begin{lstlisting}[language=Lisp]
%(progn
%  (defun fact (x)
%    (labels
%      (inner-fact (x res)
%        (if (> x 0)
%          (f (- x 1) (* res x))
%          res))
%      (inner-fact x 1))))
%->
%(progn
%  (defun fact (x res)
%    (tagbody
%      loop
%      (when (> x 0)
%        (setq res (* res x)
%              x (- x 1))
%        (go loop))
%      res))
%  (fact 5 1))
%\end{lstlisting}
%\end{figure}
%
%\item common-expression-elimination -- удаление общих подвыражений. Пример:
%\begin{figure}[H]
%\begin{lstlisting}[language=Lisp]
%(defun f (a b)
%  (let ((c (* a b 2))
%        (d (* a b 3)))
%    (+ c d)))
%->
%(defun f (a b)
%  (let ((tmp (* a b)))
%    (let ((c (* tmp 2))
%          (d (* tmp 3)))
%      (+ c d))))
%\end{lstlisting}
%\end{figure}
%
%\item continuation-passing-style -- представление программы в стиле передачи продолжений. Пример:
%\begin{figure}[H]
%\begin{lstlisting}[language=Lisp]
%(+ 1 2 3)
%->
%(+& 1 2 #'(lambda (x)
%           (+& x 3 #'(lambda (x) x))))
%\end{lstlisting}
%\end{figure}
%\end{enumerate}

\subsection{Требования к оформлению документации}

Разработка программной документации и программного изделия должна производиться согласно ГОСТ 19.102-77 и ГОСТ 34.601-90. Единая система программной документации.
